\section*{Подробные решения. Вариант 6}

\subsection*{1. Из шестнадцатеричной системы в восьмеричную}
Переводим каждую hex-цифру в четыре бита и дополняем слева нулём до групп по три:
\[
\begin{aligned}
ABC65_{16}&=\texttt{1010 1011 1100 0110 0101}_2\\
&=\texttt{010 101 011 110 001 100 101}_2\\
&=\boxed{2536145_8}.
\end{aligned}
\]

\subsection*{2. Десятичное целое в шестнадцатеричной системе}
Делим с остатком: $616=16\cdot38+8$, $38=16\cdot2+6$, $2=16\cdot0+2$. Остатки снизу вверх дают $\boxed{268_{16}}$. Проверка: $2\cdot256+6\cdot16+8=616$.

\subsection*{3. Пять шестнадцатеричных знаков дроби}
Последовательно умножаем дробный остаток на 16 и выписываем целые части:
\[
\begin{array}{rcl@{\qquad}rcl}
0.1992\cdot16&=&3.1872&0.1872\cdot16&=&2.9952,\\
0.9952\cdot16&=&15.9232&0.9232\cdot16&=&14.7712,\\
0.7712\cdot16&=&12.3392&0.3392\cdot16&=&5.4272.
\end{array}
\]
Цифры $15,14,12$ обозначаются $F,E,C$. Шестой знак 5 меньше 8, округления вверх нет: $\boxed{0.32FEC_{16}}$.

\subsection*{4. Вычитание в пятеричной системе}
Выравниваем разряды: $1234.124_5-0123.220_5$. В тысячных пишем $4-0=4$, в сотых $2-2=0$. В десятых занимаем единицу из целой части: $1+5-2=4$. Теперь целые части: $1233_5-0123_5=1110_5$.
\[
\begin{array}{r@{}l}
1234&.124_5\\
-0123&.220_5\\\hline
1110&.404_5
\end{array}
\]
Ответ $\boxed{1110.404_5}$. Проверка: разность дробных частей равна $-21/125$. После займа из целой части получаем $1-21/125=104/125=0.404_5$.

\subsection*{5. Сдвиг влево и арифметический сдвиг вправо}
Начальные коды $A=0027_{16}$ и $B=FFD0_{16}$. После логического сдвига $A$ влево на три бита получаем $0138_{16}=312_{10}$. Арифметический сдвиг $B=-48$ вправо на два бита повторяет знаковую единицу и даёт $FFF4_{16}=-12$. Разность сдвинутых значений $312-(-12)=324$, код $\boxed{0144_{16}}$.
Проверка в 16 битах: $0138_{16}-FFF4_{16}\equiv0144_{16}\pmod{2^{16}}$.

\subsection*{6. Два арифметических сдвига влево}
Код $A=-120$ равен $FF88_{16}$. Сдвиг на три бита влево даёт $FC40_{16}$, то есть $-960$. Код $B=46$ равен $002E_{16}$; после сдвига влево на четыре бита получаем $02E0_{16}=736$. Все значения входят в 16-битный знаковый диапазон. Вычитаем:
\[
-960-736=-1696,\qquad
2^{16}-1696=\boxed{F960_{16}}.
\]
Проверка кодами: $FC40_{16}-02E0_{16}=F960_{16}$.

\clearpage
\subsection*{7. Диапазон 13-битного знакового целого}
В условии не назван способ кодирования знака, поэтому минимальное значение \emph{нельзя определить однозначно}. Для 13-битного дополнительного кода:
\[
[-2^{12},\,2^{12}-1]=[-4096,\,4095].
\]
Для прямого кода (один бит знака и 12 битов модуля) и обратного кода диапазон значений равен $[-(2^{12}-1),\,2^{12}-1]=[-4095,\,4095]$. Максимум в этих трёх представлениях одинаков, минимум зависит от формата.

\subsection*{8. Функция логической схемы}
Исключающее ИЛИ вычисляет $A\oplus B$, первый И добавляет $C$, ИЛИ добавляет $D$. После НЕ результат соединяется с отдельной ветвью $A$ в последнем элементе И:
\[
\boxed{F=A\land\lnot\bigl(((A\oplus B)\land C)\lor D\bigr)}.
\]
